% Cayley, Collected Mathematical Papers, Volume X
% PDF pages 351-375 (printed pages 331-355)
% Papers: 691 (Note on Mr Monro's paper "On Flexure of Spaces.")
%         692 (Addition to the Memoir on the Transformation of Elliptic Functions)
%         693 (A Tenth Memoir on Quantics) -- opening sections through No. 374
% Scan-verified transcription of public-domain source
% Source: collectedmathema10cayluoft.pdf (Internet Archive, Cambridge 1896)

\documentclass[11pt]{article}
\usepackage[a4paper,margin=0.65in]{geometry}
\usepackage[T1]{fontenc}
\usepackage{amsmath,amssymb,amsfonts}
\usepackage{graphicx}
\usepackage{pdfpages}
\usepackage[english,shorthands=off]{babel}
\usepackage{fancyhdr}
\usepackage[bookmarks=false,colorlinks=false]{hyperref}
\usepackage[protrusion=true,expansion=false]{microtype}
\emergencystretch=3em
\tolerance=600
\providecommand{\modd}{\mathrm{mod}}
\providecommand{\Disc}{\mathrm{Disc}}
\providecommand{\canont}{}
\providecommand{\asterism}{*}
\providecommand{\Hessian}{H}
\providecommand{\tome}{}
\providecommand{\page}{p.}
\pagestyle{fancy}\fancyhf{}\fancyhead[L]{Pages 351--375}\fancyhead[R]{Cayley --- Collected Papers, Vol. X}\renewcommand{\headrulewidth}{0.4pt}
\setlength{\headheight}{15pt}

\begin{document}

% =====================================================
% Paper 691: NOTE ON MR MONRO'S PAPER "ON FLEXURE OF SPACES."
% PDF p. 351 (printed p. 331)
% =====================================================

\begin{center}
{\Large\bfseries 691.}\\[8pt]
{\large\bfseries NOTE ON MR MONRO'S PAPER ``ON FLEXURE OF SPACES.''}\\[6pt]
\textit{[From the Proceedings of the London Mathematical Society, vol.~ix.\ (1878), pp.~171, 172.\\ Read June 13, 1878.]}
\end{center}

\medskip

\noindent CONSIDER an element of surface, surrounding a point $P$; the flexure of the element may be interfered with by the continuity round $P$, and it is on this account proper to regard the element as cut or slit along a radius drawn from $P$ to the periphery of the element. This being understood, we have the well-known theorem that, considering in the neighbourhood of the origin elements of the surfaces
\[
z = \tfrac{1}{2}(ax^{2} + 2hxy + by^{2}), \quad \text{and} \quad z' = \tfrac{1}{2}(a'x'^{2} + 2h'x'y' + b'y'^{2}),
\]
these will be applicable the one on the other, provided only $ab - h^{2} = a'b' - h'^{2}$. But in connexion with Mr Monro's paper it is worth while to give the proof in detail.

It is to be shown that $z$, $z'$ denoting the above-mentioned functions of $(x, y)$ and $(x', y')$ respectively, it is possible to find (for small values) $x'$, $y'$ functions of $x, y$ such that identically
\[
dx'^{2} + dy'^{2} + dz'^{2} = dx^{2} + dy^{2} + dz^{2}.
\]

The solution is taken to be $x' = x + \xi$, $y' = y + \eta$, where $\xi$, $\eta$ denote cubic functions of $x, y$. We have then, attending only to the terms of an order not exceeding 3 in $x, y$,
\begin{align*}
& dx^{2} + dy^{2} + 2(dx\,d\xi + dy\,d\eta) + \{(a'x + h'y)\,dx + (h'x + b'y)\,dy\}^{2} \\
&\qquad = dx^{2} + dy^{2} + \{(ax + hy)\,dx + (hx + by)\,dy\}^{2},
\end{align*}
so that the terms $dx^{2} + dy^{2}$ disappear; and then writing
\[
d\xi = \tfrac{d\xi}{dx}\,dx + \tfrac{d\xi}{dy}\,dy, \quad d\eta = \tfrac{d\eta}{dx}\,dx + \tfrac{d\eta}{dy}\,dy,
\]

% ========== p. 352 (printed p. 332) ==========

\noindent the equation will be satisfied identically as regards $dx, dy$ if only
\[
2\,\tfrac{d\xi}{dx} = (ax + hy)^{2} - (a'x + h'y)^{2},
\]
\[
\tfrac{d\xi}{dy} + \tfrac{d\eta}{dx} = (ax + hy)(hx + by) - (a'x + h'y)(h'x + b'y),
\]
\[
2\,\tfrac{d\eta}{dy} = (hx + by)^{2} - (h'x + b'y)^{2}.
\]
Calling the terms on the right-hand side $2\mathfrak{A}$, $\mathfrak{B}$, $2\mathfrak{C}$ respectively, we have
\[
\frac{d^{2}\mathfrak{A}}{dy^{2}} - \frac{d^{2}\mathfrak{B}}{dx\,dy} + \frac{d^{2}\mathfrak{C}}{dx^{2}} = 0,
\]
that is,
\[
(h^{2} - h'^{2}) + (h^{2} - h'^{2}) - \{(ab - h^{2}) - (a'b' - h'^{2})\} = 0,
\]
or, what is the same thing,
\[
a'b' - h'^{2} = ab - h^{2},
\]
a relation which must exist between the constants $(a, b, h)$ and $(a', b', h')$.

It is easy to find the actual values of $\xi, \eta$; viz.\ these are
\begin{align*}
\xi &= \tfrac{1}{4}(a^{2} - a'^{2})x^{3} + \tfrac{1}{2}(ah - a'h')x^{2}y + \tfrac{1}{4}(h^{2} - h'^{2})xy^{2} + \tfrac{1}{6}(bh - b'h')y^{3},\\
\eta &= \tfrac{1}{6}(ah - a'h')x^{3} + \tfrac{1}{4}(h^{2} - h'^{2})x^{2}y + \tfrac{1}{2}(bh - b'h')xy^{2} + \tfrac{1}{4}(b^{2} - b'^{2})y^{3},
\end{align*}
or, what is the same thing, we have
\[
\xi = \tfrac{1}{4}\frac{d\Omega}{dx}, \qquad \eta = \tfrac{1}{4}\frac{d\Omega}{dy},
\]
where
\[
\Omega = (a^{2} - a'^{2})x^{4} + 4(ah - a'h')x^{3}y + 6(h^{2} - h'^{2})x^{2}y^{2} + 4(bh - b'h')xy^{3} + (b^{2} - b'^{2})y^{4}.
\]
\[
= (ax^{2} + 2hxy + by^{2})^{2} - (a'x^{2} + 2h'xy + b'y^{2})^{2} = 4(z^{2} - z'^{2}),
\]
in virtue of the relation $ab - h^{2} = a'b' - h'^{2}$. The resulting values $x' = x + \xi$, $y' = y + \eta$ are obviously the first terms of two series which, if continued, would contain higher powers of $(x, y)$.

\bigskip

% =====================================================
% Paper 692: ADDITION TO THE MEMOIR ON THE TRANSFORMATION OF ELLIPTIC FUNCTIONS.
% PDF p. 353 (printed p. 333)
% =====================================================

\begin{center}
{\Large\bfseries 692.}\\[8pt]
{\large\bfseries ADDITION TO THE MEMOIR ON THE TRANSFORMATION OF\\ ELLIPTIC FUNCTIONS.}\\[6pt]
\textit{[From the Philosophical Transactions of the Royal Society of London, vol.~CLXIX.\\
Part II.\ (1878), pp.~419--424. Received February 6,---Read March 7, 1878.]}
\end{center}

\medskip

\noindent I HAVE recently succeeded in completing a theory considered in my ``Memoir on the Transformation of Elliptic Functions,'' Phil.\ Trans., vol.\ CLXIV.\ (1874), pp.~397--456, [578],---that of the septic transformation, $n = 7$. We have here
\[
\frac{1 - y}{1 + y} = \frac{1 - x}{1 + x}\left(\frac{\alpha - \beta x + \gamma x^{2} - \delta x^{3}}{\alpha + \beta x + \gamma x^{2} + \delta x^{3}}\right)^{2},
\]
a solution of
\[
\frac{M\,dy}{\sqrt{1 - y^{2}\cdot 1 - v^{2}y^{2}}} = \frac{dx}{\sqrt{1 - x^{2}\cdot 1 - u^{2}x^{2}}},
\]
where $\dfrac{1}{M} = 1 + \dfrac{2\beta}{\alpha}$; and the ratios $\alpha : \beta : \gamma : \delta$, and the $uv$-modular equation are determined by the equations
\begin{align*}
u^{2}(2\alpha\gamma + 2\alpha\beta + \beta^{2}) &= v^{2}(\gamma^{2} + 2\gamma\delta + 2\beta\delta),\\
\gamma^{2} + 2\beta\gamma + 2\alpha\delta + 2\beta\delta &= \tfrac{1}{M^{2}}u^{2}v^{2}(2\alpha\gamma + 2\beta\gamma + 2\alpha\delta + \beta^{2}),\\
\delta^{2} + 2\gamma\delta &= \tfrac{1}{M^{2}}v^{2}u^{4}(\alpha^{2} + 2\alpha\beta);
\end{align*}
or, what is the same thing, writing $\alpha = 1$, the first equation may be replaced by
\[
\delta = \frac{v^{2}}{u},
\]
and then, $\alpha$, $\delta$ having these values, the last three equations determine $\beta$, $\gamma$ and the modular equation. If instead of $\beta$ we introduce $M$, by means of the relation

% ========== p. 354 (printed p. 334) ==========

\noindent $\dfrac{1}{M} = 1 + 2\beta$, that is, $2\beta = \dfrac{1}{M} - 1$, then the last equation gives $2\gamma = uv^{3}\left(\dfrac{1}{M} - \dfrac{u^{4}}{v^{2}}\right)$; and $\alpha$, $\beta$, $\gamma$, $\delta$ having these values, we have the residual two equations
\begin{align*}
u^{2}(2\alpha\gamma + 2\alpha\beta + \beta^{2}) &= v^{2}(\gamma^{2} + 2\gamma\delta + 2\beta\delta),\\
\gamma^{2} + 2\beta\gamma + 2\alpha\delta + \beta\delta &= u^{2}v^{2}(2\alpha\gamma + 2\beta\gamma + 2\alpha\delta + \beta^{2}),
\end{align*}
viz.\ each of these is a quadric equation in $\dfrac{1}{M}$; hence eliminating $\dfrac{1}{M}$, we have the modular equation; and also (linearly) the value of $\dfrac{1}{M}$, and thence the values of $\alpha, \beta, \gamma, \delta$ in terms of $u, v$.

Before going further it is proper to remark that, writing as above $\alpha = 1$, then if $\delta = \beta\gamma$, we have
\begin{align*}
1 - \beta x + \gamma x^{2} - \delta x^{3} &= (1 - \beta x)(1 + \gamma x^{2}),\\
1 + \beta x + \gamma x^{2} + \delta x^{3} &= (1 + \beta x)(1 + \gamma x^{2}),
\end{align*}
and the equation of transformation becomes
\[
\frac{1 - y}{1 + y} = \frac{1 - x}{1 + x}\left(\frac{1 - \beta x}{1 + \beta x}\right)^{2},
\]
viz.\ this belongs to the cubic transformation. The value of $\beta$ in the cubic transformation was taken to be $\beta = \dfrac{u^{2}}{v}$, but for the present purpose it is necessary to pay attention to an omitted double sign, and write $\beta = \pm\dfrac{u^{2}}{v}$; this being so, $\delta = \beta\gamma$, and giving to $\gamma$ the value $\mp u^{4}$, $\delta$ will have its foregoing value $-\dfrac{u^{2}}{v}$. And from the theory of the cubic equation, according as $\beta = \dfrac{u^{2}}{v}$ or $= -\dfrac{u^{2}}{v}$, the modular equation must be
\[
u^{4} - v^{4} + 2uv(1 - u^{2}v^{2}) = 0, \quad \text{or} \quad u^{4} - v^{4} - 2uv(1 - u^{2}v^{2}) = 0.
\]

We thus see \textit{\`a priori}, and it is easy to verify that the equations of the septic transformation are satisfied by the values
\begin{align*}
\alpha &= 1,\ \beta = \frac{u^{2}}{v},\ \gamma = -u^{4},\ \delta = -\frac{u^{2}}{v},\ \text{and}\ u^{4} - v^{4} + 2uv(1 - u^{2}v^{2}) = 0;\\
\alpha &= 1,\ \beta = -\frac{u^{2}}{v},\ \gamma = -u^{4},\ \delta = -\frac{u^{2}}{v},\ \text{and}\ u^{4} - v^{4} - 2uv(1 - u^{2}v^{2}) = 0;
\end{align*}
and it hence follows that in obtaining the modular equation for the septic transformation, we shall meet with the factors $u^{4} - v^{4} \pm 2uv(1 - u^{2}v^{2})$. Writing for shortness $uv = \theta$, these factors are $u^{4} - v^{4} \pm 2\theta(1 - \theta^{2})$; the factor for the proper modular equation is $u^{8} + v^{8} - \Theta$, where
\[
\Theta = 8\theta - 28\theta^{2} + 56\theta^{3} - 70\theta^{4} + 56\theta^{5} - 28\theta^{6} + 8\theta^{7},
\]

% ========== p. 355 (printed p. 335) ==========

\noindent viz.\ the equation $(1 - u^{8})(1 - v^{8}) - (1 - uv)^{8} = u^{8} + v^{8} - \Theta = 0$; and the modular equation, as obtained by the elimination from the two quadric equations, presents itself in the form
\[
(u^{4} - v^{4} + 2\theta)^{2}(u^{4} - v^{4} + 2\theta)(u^{4} + v^{4} - \Theta) = 0.
\]

Proceeding to the investigation, we substitute the values
\[
\alpha = 1,\ \beta = \tfrac{1}{2}\left(\tfrac{1}{M} - 1\right),\ \gamma = \tfrac{1}{2}uv^{3}\left(\tfrac{1}{M} - \tfrac{u^{4}}{v^{2}}\right),\ \delta = \tfrac{v^{2}}{u},
\]
in the residual two equations, which thus become
\begin{align*}
\frac{1}{M^{2}}(1 - v^{4}) &+ \frac{2}{M}(1 - uv)^{4}(1 + uv) \\
&\qquad + \{(1 - v^{8}) - 4(1 - uv)(1 + uv)^{3}\} = 0,\\[4pt]
\frac{1}{M^{2}}\{-u^{4}v^{2}(1 - uv)^{4}(1 + uv)\} &+ \frac{2}{M}\{u^{2}v^{2}(1 - v^{4}) + \tfrac{u^{2}}{v}(1 + u^{2}v^{2})(u^{4} - v^{4})\}\\
&\qquad + \{u^{8} + 6\tfrac{u^{4}}{v^{2}}(1 - u^{2}v^{2}) - \tfrac{u^{4}}{v^{2}}\}(1 - v^{4}) = 0,
\end{align*}
the first of which is given p.~432 of the ``Memoir,'' [Coll.\ Math.\ Papers, vol.\ ix., p.~150]. Calling them
\[
(a, b, c\,\$\,\tfrac{1}{M}, 1)^{2} = 0, \quad (a', b', c'\,\$\,\tfrac{1}{M}, 1)^{2} = 0,
\]
we have
\[
\frac{-1}{bc' - b'c} : \frac{2}{ca' - c'a} : \frac{1}{ab' - a'b},
\]
and the result of the elimination therefore is
\[
(ca' - c'a)^{2} - 4(bc' - b'c)(ab' - a'b) = 0.
\]

Write as before $uv = \theta$. In forming the expressions $ca' - c'a$, \&c., to avoid fractions we must in the first instance introduce the factor $v^{2}$: thus
\begin{align*}
v^{2}(ca' - c'a) &= v^{2}\{v^{2}(1 - u^{8}) - 4(1 - \theta)(1 + \theta)^{3}\}\cdot[-u^{4}v^{2}(1 - \theta)^{4}(1 + \theta)]\\
&\quad - \{u^{8} + 6u^{4}\theta^{2}(1 - \theta^{2}) - v^{4}\theta^{2}\}\{1 - v^{4}\},\\
&= -u^{4}v^{4}(1 + \theta)(1 - \theta)^{4}\{v^{2}(-3 + 4\theta) + u^{4}(-4\theta + 3\theta^{2})\}\\
&\quad - \{u^{8} + 6u^{4}(\theta - \theta^{3}) - v^{4}\theta^{2}\}(1 - v^{4})
\end{align*}
but instead of $\theta^{2}v^{2}$ writing $v^{6}$, the expression on the right-hand side becomes divisible by $v^{4}$; and we find
\begin{align*}
\frac{1}{v^{4}}(ca' - c'a) &= -(1 + \theta)(1 - \theta)^{4}\{v^{4}(-3 + 4\theta) + u^{4}(-4\theta + 3\theta^{2})\}\\
&\quad - \{u^{8} + 6u^{4}(\theta - \theta^{3}) - v^{4}\}(1 - v^{4}),
\end{align*}

% ========== p. 356 (printed p. 336) ==========

\noindent and thence
\begin{align*}
\frac{1}{v^{4}}(ca' - c'a) &= u^{8} + u^{4}(6\theta - 10\theta^{2} + 11\theta^{3} - 6\theta^{4} - 8\theta^{5} + 10\theta^{6} - 4\theta^{7})\\
&\quad + v^{4}(-4 + 10\theta - 8\theta^{2} - 6\theta^{3} + 11\theta^{4} - 10\theta^{5} + 6\theta^{6}) + v^{8}.
\end{align*}
Similarly we have
\begin{align*}
\frac{1}{v^{4}}(bc' - b'c) &= u^{8}(5 - 5\theta + 4\theta^{2} - 5\theta^{3} + 2\theta^{4}) + u^{4}(9\theta - 16\theta^{2} + \theta^{3} + 10\theta^{4} + \theta^{5} - 16\theta^{6} + 9\theta^{7})\\
&\quad + v^{4}(2 - 5\theta + 4\theta^{2} - 5\theta^{3} + 5\theta^{4}),\\[4pt]
\frac{1}{v^{4}}(ab' - a'b) &= u^{8}(\theta + \theta^{2} - \theta^{4}) + u^{4}(2 - 5\theta + 4\theta^{2} + 3\theta^{3} - 10\theta^{4} + 3\theta^{5} + 4\theta^{6} - 5\theta^{7} + 2\theta^{8})\\
&\quad + v^{4}(-1 + \theta + \theta^{2});
\end{align*}
say these values are
\[
u^{8} + p u^{4} + q v^{4} + v^{8},\quad \lambda u^{8} + \mu u^{4} + \nu v^{4},\quad \rho u^{8} + \sigma u^{4} + \tau v^{4}.
\]
The required equation is thus
\[
(u^{8} + p u^{4} + q v^{4} + v^{8})^{2} - 4(\lambda u^{8} + \mu u^{4} + \nu v^{4})(\rho u^{8} + \sigma u^{4} + \tau v^{4}),
\]
viz.\ the function is
\begin{align*}
& u^{16}\\
&+ u^{12}(2p - 4\lambda\rho)\\
&+ u^{8}(2q\theta^{4} + p^{2} - 4\lambda\sigma\theta^{4} - 4\mu\rho)\\
&+ (2p\theta^{4} + 2pq\theta^{4} - 4\lambda\tau\theta^{8} - 4\mu\sigma\theta^{4} - 4\nu\rho\theta^{4})\\
&+ u^{4}(2p\theta^{4} + q^{2}\theta^{4} - 4\mu\tau\theta^{8} - 4\nu\sigma\theta^{4})\\
&+ v^{12}(2q - 4\nu\tau)\\
&+ v^{16},
\end{align*}
or say it is
\[
= (1, b, c, d, e, f\,\$\,u^{16}, u^{12}, u^{8}, 1, v^{8}, v^{12}, v^{16}).
\]

Supposing that this has a factor $u^{8} + v^{8} - \Theta$, the form is
\[
(u^{16} + Bu^{8} + Cu^{4}v^{4} + Dv^{8} + v^{16})(u^{8} + v^{8} - \Theta)
\]
and comparing coefficients we have
\begin{align*}
B - \Theta &= b,\\
C - \Theta B + \Theta^{2} &= c,\\
D\Theta^{2} - \Theta C + B\Theta^{3} &= d,\\
\Theta^{2} - \Theta D + C &= e,\\
-\Theta + D &= f,
\end{align*}
where $\Theta$ has the before-mentioned value
\[
\Theta = (8, -28, +56, -70, +56, -28, +8\,\$\,\theta, \theta^{2}, \theta^{3}, \theta^{4}, \theta^{5}, \theta^{6}, \theta^{7}).
\]
From the first, second, and fifth equations, $B = b + \Theta$, $D = f + \Theta$, $C = c + \Theta b + \Theta^{2}$; and

% ========== p. 357 (printed p. 337) ==========

\noindent the third and fourth equations should then be verified identically. Writing down the coefficients of the different powers of $\theta$, we find
\begin{align*}
2p &= 0 + 12 \quad 0 - 20 + 22 - 12 - 16 + 20 - 8\quad (\theta, \ldots, \theta^{7})\\
-4\lambda\rho &= -20 + 20 - 36 + 60 - 44 + 36 - 28 + 8\\
b &= 0 - 8 + 20 - 56 + 82 - 56 + 20 - 8\\
e &= 0 + 8 - 28 + 56 - 70 + 56 - 28 + 8\\
B &= 0 - 8 \quad + 12 \quad 0 - 8
\end{align*}
that is,
\[
B = -8\theta + 12\theta^{2} - 8\theta^{3};
\]
and in precisely the same way the fifth equation gives the corresponding value of $D$.

We find similarly $C$ from the second equation: writing down first the coefficients of $p^{2}$, $2q\theta^{4}$, $-4\lambda\sigma\theta^{4}$, and $4\mu\rho$, the sum of these gives the coefficients of $c$; and then writing underneath these the coefficients of $B\Theta$ and of $-\Theta^{2}$, the final sum gives the coefficients of $C$: the coefficients of each line belong to $(\theta^{0}, \theta^{1}, \ldots, \theta^{12})$.
\[
\resizebox{\textwidth}{!}{$\displaystyle
\begin{array}{r}
36\quad 0 - 120 + 132 + 28 - 316 + 361 - 20 - 340 + 396 - 144 - 112 + 164 - 80 + 16\\
-8 + 20 - 16 - 12 + 22 - 20 \quad 0 + 12 \\
-40 + 140 - 212 + 140 + 80 - 188 + 168 - 92 - 64 + 176 - 164 + 80 - 16\\
-36 + 64 - 40 + 60 - 72 + 28 \quad 0 + 68 - 100 + 36\\
0 + 64 - 208 + 352 - 272 - 160 + 463 - 160 - 272 + 352 - 208 + 64\\
-64 + 224 - 352 + 224 + 160 - 392 + 160 + 224 - 352 + 224 - 64\\
\hline
0\quad 0\quad 0\quad 0 + 16 \quad 0 - 48 \quad 0 + 70 \quad 0 - 48 \quad 0 + 16 \quad 0,
\end{array}$}
\]
that is,
\[
C = 16\theta^{4} - 48\theta^{5} + 70\theta^{6} - 48\theta^{7} + 16\theta^{8};
\]
and in precisely the same way this value of $C$ would be found from the fourth equation. There remains to be verified only the fourth equation $(D + B)\Theta^{2} - \Theta C = d$, that is,
\[
2f\Theta^{2}(-8\theta + 12\theta^{2} - 8\theta^{3}) - \Theta C = (2p - 4\lambda\tau)\theta^{8} + (2pq - 4\mu\sigma - 4\nu\rho)\theta^{4},
\]
and this can be effected without difficulty.

The factor of the modular equation thus is
\[
u^{16} + v^{16} + (-8\theta + 12\theta^{2} - 8\theta^{3})(u^{8} + v^{8}) + 16\theta^{4} - 48\theta^{5} + 70\theta^{6} - 48\theta^{7} + 16\theta^{8}
\]

% ========== p. 358 (printed p. 338) ==========

\noindent viz.\ this is
\begin{align*}
(u^{8} + v^{8})^{2} &+ (-4\theta + 6\theta^{2} - 4\theta^{3}) \cdot 2(u^{8} + v^{8}) + 16\theta^{2} - 48\theta^{3} + 68\theta^{4} - 48\theta^{5} + 16\theta^{6}\\
&= \{u^{8} + v^{8} - 4\theta + 6\theta^{2} - 4\theta^{3}\}^{2}\\
&= \{(u^{4} - v^{4})^{2} - 4\theta(1 - \theta^{2})\}^{2},
\end{align*}
that is,
\[
\{u^{4} - v^{4} - 2\theta(1 - \theta^{2})\}^{2}\{u^{4} - v^{4} + 2\theta(1 - \theta^{2})\}^{2};
\]
or the modular equation is
\[
\{u^{4} - v^{4} - 2\theta(1 - \theta^{2})\}^{2}\{u^{4} - v^{4} + 2\theta(1 - \theta^{2})\}^{2}(u^{8} + v^{8} - \Theta) = 0;
\]
viz.\ the first and second factors belong to the cubic transformation; and we have for the proper modular equation in the septic transformation $u^{8} + v^{8} - \Theta = 0$, or what is the same thing $(1 - u^{8})(1 - v^{8}) - (1 - \theta)^{8} = 0$, that is, $(1 - u^{8})(1 - v^{8}) - (1 - uv)^{8} = 0$, the known result; or, as it may also be written,
\[
\{(1 - u^{2})(1 - v^{2})\}^{4} = (1 - uv)^{8}.
\]

The value of $M$ is given by the foregoing relations
\[
\tfrac{-1}{\lambda u^{8} + \mu u^{4} + \nu v^{4}} : \tfrac{2}{u^{8} + pu^{4} + qv^{4} + v^{8}} : \tfrac{1}{\rho u^{8} + \sigma u^{4} + \tau v^{4}};
\]
but these can be, by virtue of the proper modular equation $u^{8} + v^{8} - \Theta = 0$, reduced into the form

viz.\ the equality of these two sets of ratios depends upon the following identities,
\begin{align*}
&(-\lambda u^{8} + \nu v^{8})(u^{8} + pu^{4} + qv^{4} + v^{8}) + (\theta - 2\theta^{2} + 2\theta^{3} - \theta^{4})(\rho u^{8} + \sigma u^{4} + \tau v^{4})\\
&\qquad = \{-\tfrac{1}{2}u^{8} + (1 - \theta)(-4 - \theta + 5\theta^{2} - \theta^{3} - 4\theta^{4})v^{4} + v^{8}\}(u^{8} - \Theta + v^{8}),\\[4pt]
&(\theta - u^{4})(\rho u^{8} + \sigma u^{4} + \tau v^{4}) - (\theta - v^{4})(\lambda u^{8} + \mu u^{4} + \nu v^{4})\\
&\qquad = \{(2\theta + 5\theta^{2} + 3\theta^{3} - 2\theta^{4} - 2\theta^{5})u^{4} + (2 + 2\theta - 3\theta^{2} - 5\theta^{3} - 2\theta^{4})v^{4}\}(u^{8} - \Theta + v^{8}),\\[4pt]
&-2(\theta - 2\theta^{2} + 2\theta^{3} - \theta^{4})(\lambda u^{8} + \mu u^{4} + \nu v^{4}) + (u^{8} - \theta)(u^{8} + pu^{4} + qv^{4} + v^{8})\\
&\qquad = \{u^{8} + \theta(1 - \theta)(3 + 5\theta + 3\theta^{2})u^{4} - \theta^{4}v^{4}\}(u^{8} - \Theta + v^{8}),
\end{align*}
which can be verified without difficulty: from the last-mentioned system of values, replacing $\theta$ by its value $uv$, we then have
\[
\frac{-1}{\tfrac{1}{2}u(v - u^{7})} : \frac{2}{1 - uv(1 - uv)(1 - uv + u^{2}v^{2})} : \frac{1}{-v(u - v^{7})},
\]
which agree with the values given p.~432 of the ``Memoir''; and the analytical theory is thus completed.

\bigskip

% =====================================================
% Paper 693: A TENTH MEMOIR ON QUANTICS.
% PDF p. 359 (printed p. 339)
% =====================================================

\begin{center}
{\Large\bfseries 693.}\\[8pt]
{\large\bfseries A TENTH MEMOIR ON QUANTICS.}\\[6pt]
\textit{[From the Philosophical Transactions of the Royal Society of London, vol.~CLXIX., Part II.\\ (1878), pp.~603--661. Received June 12,---Read June 20, 1878.]}
\end{center}

\medskip

\noindent The present Memoir, which relates to the binary quintic $(*\,\$\,x, y)^{5}$, has been in hand for a considerable time: the chief subject-matter was intended to be the theory of a canonical form which was discovered by myself and is briefly noticed in Salmon's \textit{Higher Algebra}, 3rd Ed.\ (1876), pp.~217, 218; writing $a, b, c, d, e, f, g, \ldots, u, v, w$ to denote the 23 covariants of the quintic, then $a, b, c, d, f$ are connected by the relation
\[
f^{2} = -a^{3}d + a^{2}bc - 4c^{3};
\]
and the form contains these covariants thus connected together, and also $e$; it, in fact, is
\[
(1, 0, c, f, a^{2}b - 3c^{2}, a^{2}e - 2cf\,\$\,x, y)^{5}.
\]

But the whole plan of the Memoir was changed by Sylvester's discovery of what I term the Numerical Generating Function (N.G.F.) of the covariants of the quintic, and my own subsequent establishment of the Real Generating Function (R.G.F.) of the same covariants. The effect of this was to enable me to establish for any given degree in the coefficients and order in the variables, or as it is convenient to express it, for any given deg-order whatever, a selected system of powers and products of the covariants, say a system of ``segregates'': these are asyzygetic, that is, not connected together by any linear equation with numerical coefficients; and they are also such that every other combination of covariants of the same deg-order, say every ``congregate'' of the same deg-order, can be expressed (and that, obviously, in one way only) as a linear function, with numerical coefficients, of the segregates of that deg-order. The number of congregates of a given deg-order is precisely equal to the number of the independent syzygies of the same deg-order, so that these syzygies give in effect the congregates in terms of the segregates: and the proper form in which to exhibit the

% ========== p. 360 (printed p. 340) ==========

\noindent syzygies is thus to make each of them give a single congregate in terms of the segregates: viz.\ the left-hand side can always be taken to be a monomial congregate $a^{\alpha}b^{\beta}\ldots$, or, to avoid fractions, a numerical multiple of such form; and the right-hand side will then be a linear function, with numerical coefficients, of the segregates of the same deg-order. Supposing such a system of syzygies obtained for a given deg-order, any covariant function (rational and integral function of covariants) is at once expressible as a linear function of the segregates of that deg-order: it is, in fact, only necessary to substitute therein for every monomial congregate its value as a linear function of the segregates. Using the word covariant in its most general sense, the conclusion thus is that every covariant can be expressed, and that in one way only, as a linear function of segregates, or say in the segregate form.

Reverting to the theory of the canonical form, and attending to the relation
\[
f^{2} = -a^{3}d + a^{2}bc - 4c^{3},
\]
it thereby appears that every covariant multiplied by a power of the quintic itself $a$, can be expressed, and that in one way only, as a rational and integral function of the covariants $a, b, c, d, e, f$, linear as regards $f$: say every covariant multiplied by a power of $a$ can be expressed, and that in one way only, in the ``standard'' form as an illustration, take
\[
a^{2}h = 6acd + 4bc^{2} + ef.
\]
Conversely, an expression of the standard form, that is, a rational and integral function of $a, b, c, d, e, f$, linear as regards $f$, not explicitly divisible by $a$, may very well be really divisible by a power of $a$ (the expression of the quotient of course containing one or more of the higher covariants $g, h$, \&c.), and we say that in this case the expression is divisible, and has for its divided form the quotient expressed as a rational and integral function of covariants. Observe that in general the divided form is not perfectly definite, only becoming so when expressed in the before-mentioned segregate form, and that this further reduction ought to be made. There is occasion, however, to consider these divided forms, whether or not thus further reduced; and moreover it sometimes happens that the non-segregate form presents itself, or can be expressed, with integer numerical coefficients, while the coefficients of the corresponding segregate form are fractional.

The canonical form is peculiarly convenient for obtaining the expressions of the several derivatives (Gordan's \textit{Uebereinanderschiebungen}) $(a, b)^{1}$, $(a, b)^{2}$, \&c., (or as I propose to write them $ab1$, $ab2$, \&c.), which can be formed with two covariants, the same or different, as rational and integral functions of the several covariants. It will be recollected that in Gordan's theory these derivatives are used in order to establish the system of the 23 covariants: but it seems preferable to have the system of covariants, and by means of them to obtain the theory of the derivatives.

I mention at the end of the Memoir two expressions (one or both of them due to Sylvester) for the N.G.F.\ of a binary sextic.

The several points above adverted to are considered in the Memoir; the paragraphs are numbered consecutively with those of the former Memoirs upon Quantics.

% ========== p. 361 (printed p. 341) ==========

\begin{center}
\textit{The Numerical and Real Generating Functions.} Art.\ Nos.\ 366 to 374,\\
and Table No.\ 96.
\end{center}

\medskip

\noindent 366. I have, in my Ninth Memoir (1871) [462], given what may be called the Numerical Generating Function (N.G.F.) of the covariants of a quartic; this was
\[
A(x) = \frac{1 - a^{6}x^{12}}{1 - ax^{4}\cdot 1 - a^{2}x^{4}\cdot 1 - a^{2}\cdot 1 - a^{3}\cdot 1 - a^{3}x^{6}},
\]
the meaning being that the number of asyzygetic covariants $a^{\theta}x^{\mu}$, of the degree $\theta$ in the coefficients and order $\mu$ in the variables, or say of the deg-order $\theta.\mu$, is equal to the coefficient of $a^{\theta}x^{\mu}$ in the development of this function. And I remarked that the formula indicated that the covariants were made up of $(ax^{4}, a^{2}x^{4}, a^{2}, a^{3}, a^{3}x^{6})$, the quartic itself, the Hessian, the quadrinvariant, the cubinvariant, and the cubicovariant, these being connected by a syzygy $a^{6}x^{12}$ of the degree 6 and order 12. Calling these covariants $a, b, c, d, e$, so that these italic small letters stand for covariants,
\begin{center}
\begin{tabular}{ll}
deg-order. & \\
1.\,4 & $a$,\\
2.\,0 & $b$,\\
2.\,4 & $c$,\\
3.\,0 & $d$,\\
3.\,6 & $e$,\\
\end{tabular}
\end{center}
then it is natural to consider what may be called the Real Generating Function (R.G.F.): this is
\[
\frac{1 - e^{2}}{1 - a\cdot 1 - b\cdot 1 - c\cdot 1 - d\cdot 1 - e},
\]
the development of this contains, as it is easy to see, only terms of the form $a^{\alpha}b^{\beta}c^{\gamma}d^{\delta}$ and $a^{\alpha}b^{\beta}c^{\gamma}d^{\delta}e$, each with the coefficient $+1$, so that the number of terms of a given deg-order $\theta.\mu$ is equal to the coefficient of $a^{\theta}x^{\mu}$ in the first-mentioned function: and these terms of a given deg-order represent the asyzygetic covariants of that deg-order: any other covariant of the same deg-order is expressible as a linear function of them.

For instance, deg-order 6.12, the terms of the R.G.F.\ are $a^{3}d$, $a^{2}bc$, $c^{3}$; there is one more term $e^{2}$ of the same deg-order; hence $e^{2}$ must be a linear function of these and in fact
\[
e^{2} = -a^{3}d + a^{2}bc - 4c^{3},
\]
viz.\ this is the equation
\[
\Phi^{2} = -U^{3}J + U^{2}IH - 4H^{3}.
\]

% ========== p. 362 (printed p. 342) ==========

\noindent 367. Sylvester obtained an expression for the N.G.F.\ of the quintic: this is
\[
\resizebox{\textwidth}{!}{$\displaystyle
\frac{
\begin{aligned}
& a^{0}\cdot 1\\
& + a^{3}\cdot x^{3} + x^{5} + x^{9}\\
& + a^{4}\cdot x^{4} + x^{6}\\
& + a^{5}\cdot x + x^{3} + x^{7}\\
& + a^{6}\cdot x^{2} + x^{4}\\
& + a^{7}\cdot x + x^{3} - x^{5} - x^{7}\\
& + a^{8}\cdot x^{2} - x^{4}\\
& + a^{9}\cdot x + a^{3}\cdot x^{3}\\
& + a^{10}\cdot x^{2} + x^{4} - x^{6}\\
& + a^{11}\cdot x - x^{3} - x^{5}\\
& + a^{12}\cdot x^{0} - x^{2} - x^{6} - x^{10}\\
& + a^{13}\cdot x - x^{3} - x^{7}\\
& + a^{14}\cdot x^{2} - x^{4} - x^{6}\\
& + a^{15}\cdot -x^{3} - x^{5}\\
& + a^{16}\cdot -x^{2} - x^{4} - x^{8}\\
& + a^{17}\cdot -x - x^{3}\\
& + a^{18}\cdot -x^{4} - x^{6} - x^{10}\\
& + a^{19}\cdot x^{5} - x^{7}\\
& + a^{23}\cdot -x^{3} - x^{5} - x^{9}
\end{aligned}
}{1 - ax^{5}\cdot 1 - a^{2}x^{2}\cdot 1 - a^{2}x^{6}\cdot 1 - a^{3}x^{3}\cdot 1 - a^{4}\cdot 1 - a^{8}}$}
\]
viz.\ expanding this function in ascending powers of $a, x$, then, if a term is $N a^{\theta}x^{\mu}$, this means that there are precisely $N$ asyzygetic covariants of the deg-order $\theta.\mu$.

368. It is known that the number of the irreducible covariants of the binary quintic is $=23$; representing these by the letters $a, b, c, d, e, f, g, h, i, j, k, l, m, n, o, p, q, r, s, t, u, v, w$ ($a$ the quintic itself), the deg-orders of these, and the references$^{*}$ to the tables which give them are

\medskip

\noindent\textit{[* See also the paper, 143, in the second volume of this collection.]}

% ========== p. 363 (printed p. 343) ==========

\begin{center}
\begin{tabular}{lll}
\textbf{Deg-order.} & & \textbf{Tab.\ Mem.}\\
1\,.\,5 & $a$ & 13 \quad 2\\
2\,.\,2 & $b$ & 14\\
,,\,.\,6 & $c$ & 15\\
3\,.\,3 & $d$ & 16\\
,,\,.\,5 & $e$ & 17\\
,,\,.\,9 & $f$ & 18\\
4\,.\,0 & $g$ & 19\\
,,\,.\,4 & $h$ & 20\\
,,\,.\,6 & $i$ & 21\\
5\,.\,1 & $j$ & 22\\
,,\,.\,3 & $k$ & 23\\
,,\,.\,7 & $l$ & 24\\
6\,.\,2 & $m$ & 83 \quad 8\\
,,\,.\,4 & $n$ & 84\\
7\,.\,1 & $o$ & 90$^{*}$ \quad 9\\
,,\,.\,5 & $p$ & 91 \quad 2\\
8\,.\,0 & $q$ & 25 \quad 2 \ [See also paper 143]\\
,,\,.\,2 & $r$ & 92 \quad 9\\
9\,.\,3 & $s\dagger$ & 94 \quad 9\\
11\,.\,1 & $t$ & 29 \quad 3\\
12\,.\,0 & $u$ & 95 \quad 9\\
13\,.\,1 & $v$ & 29a \quad 5.\\
18\,.\,0 & $w$ & \\
\end{tabular}
\end{center}

\noindent Starting from the foregoing expression of the N.G.F.\ of the quintic, we can, instead of each term $a^{\theta}x^{\mu}$, introduce a covariant or product of covariants of the proper deg-order $\theta.\mu$: the mode of doing this depends of course on the different admissible partitions of $\theta, \mu$, and it is for some of the terms very indeterminate: for instance, $a^{4}x^{6}$ is $ai$, $hj$, or $ce$. I found it possible to perform the whole process so as to satisfy a condition which will be presently referred to; and I found

\medskip

\noindent {\footnotesize [* See vol.\ VII.\ of this collection, p.\ 348.]\\
$\dagger$ See end of Memoir. The $S$ of Table 93 has the value $-96(D, M) + 16BO - 7GA'$, but it is better to use the simple value $-(D, M)$; and the $S$ of the present Memoir has this value, say $S = -(d, m)$.}

% ========== p. 364 (printed p. 344) ==========

\begin{center}
\textit{R.G.F.\ of quintic:}
\end{center}
\[
\resizebox{\textwidth}{!}{$\displaystyle
\begin{array}{ll}
1. & \text{Deg-orders.}\\
& 0.0\text{--}10.10\\
+d \cdot 1 - ag^{2} & 3.3\text{--}12.\,8\\
& 3.5\text{--}\,7.\,9\\
+e \cdot 1 - b^{2} & 3.9\text{--}\,5.11\\
& 4.4\text{--}13.\,9\\
+f \cdot 1 - b & 4.6\text{--}12.10\\
+h \cdot 1 - ag^{2} & 5.1\text{--}14.\,6\\
& 5.3\text{--}\,9.\,7\\
+i \cdot 1 - b^{2}g & 5.7\text{--}11.\,9\\
+j \cdot 1 - ag^{2} & 6.2\text{--}15.\,7\\
& 6.4\text{--}14.\,8\\
+k \cdot 1 - b^{2} & 7.1\text{--}13.\,7\\
& 7.5\text{--}15.\,9\\
+l \cdot 1 - bg & 8.2\text{--}16.\,6\\
& 8.4\text{--}17.\,9\\
+m \cdot 1 - ag^{2} & 9.3\text{--}16.10\\
+n \cdot 1 - b^{2}g & 9.5\text{--}18.10\\
+o \cdot 1 - b^{2} & 10.2\text{--}19.\,7\\
& 10.4\text{--}18.\,8\\
+p \cdot 1 - b^{2}g & 11.1\text{--}17.\,7\\
+r \cdot 1 - \tfrac{1}{6} & 11.3\text{--}20.\,8\\
+dj \cdot 1 - ag^{2} & 12.2\text{--}18.\,4\\
+s \cdot 1 - abg & 13.1\text{--}23.11\\
+hj \cdot 1 - ag^{2} & 14.4\text{--}20.\,6\\
+j^{2} \cdot 1 - ag^{2} & 16.2\text{--}20.\,2\\
+jk \cdot 1 - b^{2}g & 18.0\text{--}19.\,5\\
+t \cdot 1 - b^{2} & \\
+jm \cdot 1 - ag^{2} & \\
+jo \cdot 1 - bg & \\
+v \cdot 1 - b^{2} & \\
+js \cdot 1 - bg & \\
+jt \cdot 1 - g & \\
+w \cdot 1 - a & \\[4pt]
\multicolumn{2}{l}{1 - a\cdot 1 - b\cdot 1 - c\cdot 1 - g\cdot 1 - q\cdot 1 - u,}
\end{array}$}
\]
where observe that each negative term of the numerator is equal to a positive term multiplied by a power or product of terms $a, b, g$, contained in the denominator: this is the condition above referred to. The expansion thus consists only of terms each with the coefficient $+1$. For instance, a part of the function is
\[
\frac{s(1 - abg)}{1 - a\cdot 1 - b\cdot 1 - c\cdot 1 - g\cdot 1 - q\cdot 1 - u} = \frac{s}{1 - c\cdot 1 - q\cdot 1 - u}\cdot\frac{1 - abg}{1 - a\cdot 1 - b\cdot 1 - g},
\]

% ========== p. 365 (printed p. 345) ==========

\noindent where the first factor is the entire series of terms $s c^{\alpha} q^{\beta} u^{\gamma}$, and the second factor is the series of terms $a^{\alpha}b^{\beta}g^{\gamma}$ omitting only those terms which are divisible by $abg$: and in the product of the two factors the terms are all distinct, so that the coefficients are still each $=1$. The same thing is true for every other pair of numerator terms and since the terms arising from each such pair are distinct from each other, in the expansion of the entire function the coefficients are each $=+1$. Hence (as in the case of the quartic) for any given deg-order, the terms in the expansion of the R.G.F.\ may be taken for the asyzygetic covariants of that deg-order; and if there are any other terms of the same deg-order, each of these must be a linear function, with numerical coefficients, of these asyzygetic covariants: thus deg-order 6\,.\,14, the expansion contains only the terms $a^{2}h$, $acd$, $bc^{2}$; there is besides a term of the same deg-order, $ef$, which is not a term of the expansion, and hence $ef$ must be a linear function of $a^{2}h$, $acd$, $bc^{2}$; we in fact have $ef = a^{2}h - 6acd - 4bc^{2}$.

The terms in the expansion of the R.G.F.\ may be called ``segregates,'' and the terms not in the expansion ``congregates''; the theorem thus is: every congregate is a linear function, with determinate numerical coefficients, of the segregates of the same deg-order.

369. I stop to remark that the numerator of the R.G.F.\ may be written in the more compendious form
\begin{align*}
& (1 - b^{2})(1 - v) + (1 - b^{2})(o + t) + (1 - b)(e + k) + (1 - b)f\\
& \quad + (1 - ag^{2})(d + h + j + m + dj + hj + j^{2} + jm)\\
& \quad + (1 - bg)(l + jo + js)\\
& \quad + (1 - b^{2}g)(n + p + i + jk)\\
& \quad + (1 - abg)s\\
& \quad + (1 - g)jt\\
& \quad + (1 - a)w;
\end{align*}
but the first-mentioned form is, I think, the more convenient one.

370. It is to be noticed that the positive terms of the numerator are unity, the seventeen covariants $d, e, f, h, i, j, k, l, m, n, o, p, r, s, t, v, w$, and the products of $j$ by $(d, h, j, k, m, o, s, t)$, where $j^{2}$ is reckoned as a product; in all, 26 terms. Disregarding the negative terms of the numerator the expansion would consist of these 26 terms, each multiplied by every combination whatever $a^{\alpha}b^{\beta}c^{\gamma}g^{\delta}q^{\varepsilon}u^{\zeta}$ of the denominator terms $a, b, c, g, q, u$ (which for this reason might be called ``reiterative''): the effect of the negative terms of the numerator is to remove from the expansion certain of the terms in question, thereby diminishing the number of the segregates: thus as regards the terms belonging to unity, any one of these which contains the factor $b^{2}$ is not a segregate but a congregate: and so as regards the terms belonging to $d$, any one of these which contains the factor $ag^{2}$ is a congregate: and the like in other cases.

For a given deg-order we have a certain number of segregates and a certain number of congregates: and the number of independent syzygies of that deg-order is

% ========== p. 366 (printed p. 346) ==========

\noindent precisely equal to the number of congregates: viz.\ each such syzygy may be regarded as giving a congregate in terms of the segregates: we have on the left-hand side a congregate, or, to avoid fractions, a numerical multiple of the congregate, and on the right-hand side a linear function, with numerical coefficients, of the segregates.

371. The syzygy is irreducible or reducible; and in the latter case it is, or is not, simply divisible: viz.\ if the congregate on the left-hand side contains any congregate factor (the other factor being literal), then the syzygy is reducible: it is, in fact, obtainable from the syzygy (of a lower deg-order) which gives the value of such congregate factor. But there are here two cases; multiplying the lower syzygy by the proper factor, the right-hand side may still contain segregates only, and then no further step is required: the original syzygy is nothing else than this lower syzygy, each side multiplied by the factor in question, and it is accordingly said to be simply divisible (S.D.). But contrariwise, the right-hand side, as multiplied, may contain congregates which have to be replaced by their values in terms of the segregates of the same deg-order: the resulting expression is then no longer explicitly divisible by the introduced factor: and the original syzygy, although arising as above from a lower syzygy, is not this lower syzygy each side multiplied by a factor: viz.\ it is in this case not simply divisible.

For example (see the subsequent Table No.\ 96, under the indicated deg-orders) (6\,.\,6), from the syzygy
\[
9d^{2} = aj - b^{2} + 2bh - cg,
\]
we deduce (7\,.\,11) the syzygy
\[
9da^{2} = a^{2}j - ab^{2} + 2abh - acg,
\]
which (all the terms on the right-hand being segregates) requires no further reduction it is a reducible and simply divisible syzygy. But we have (6\,.\,8) a syzygy giving $de$, and also (6\,.\,10) a syzygy giving $e^{2}$; multiplying the former of these by $e$ or the latter of them by $d$, we obtain values of $de^{2}$, but in each case the right-hand sides contain terms which are not segregates, and have thus to be further reduced; the final formula (9\,.\,13) is
\[
3a^{2}de^{2} = -3a^{2}d^{2}f + 4a^{2}bj + 8abh^{2} - 4abcg + 12b^{2}cd,
\]
which is not divisible by any factor: the syzygy is thus reducible, but not simply divisible.

A syzygy, which is not in the sense explained reducible, is said to be irreducible.

372. The number of irreducible syzygies is obviously finite: it has, however, the large value 179 as appears from the annexed diagram, showing the congregates determined by these several syzygies, and the deg-orders of the syzygies:

% ========== p. 367 (printed p. 347) ==========

\clearpage

\begin{center}
\begingroup
\scriptsize
\setlength{\tabcolsep}{2.2pt}
\renewcommand{\arraystretch}{1.15}
\newcommand{\dg}[1]{\ensuremath{#1}}
\newcommand{\fg}[2]{\shortstack{\ensuremath{#1}\\[-1pt]\ensuremath{#2}}}
\resizebox{\textwidth}{!}{%
\begin{tabular}{@{}r c|*{19}{c}|c@{}}
 & & \dg{1} & \dg{d} & \dg{e} & \dg{f} & \dg{h} & \dg{i} & \dg{j} & \dg{k} & \dg{l} & \dg{m} & \dg{n} & \dg{o} & \dg{p} & \dg{r} & \dg{s} & \dg{t} & \dg{v} & \dg{w} & \dg{j^{2}} & \\ \hline
\dg{0.0} & \dg{1} &
  \fg{b^{5}}{10.10} & \fg{ag^{2}}{12.8} & \fg{b^{2}}{7.9} & \fg{b}{5.11} & \fg{ag^{2}}{13.9} & \fg{b^{2}g}{12.10} & \fg{ag^{2}}{14.6} & \fg{b^{2}}{9.7} & \fg{bg}{11.9} & \fg{ag^{2}}{15.7} & \fg{b^{2}g}{14.8} & \fg{b^{3}}{13.7} & \fg{b^{2}g}{15.9} & \fg{b^{2}g}{16.6} & \fg{abg}{16.10} & \fg{b^{3}}{17.7} & \fg{b^{5}}{23.11} & \fg{a}{19.5} & & \dg{1} \\
\dg{3.3} & \dg{d} &
  & \dg{6.6} & \dg{6.8} & \dg{6.12} & \dg{7.7} & \dg{7.9} & \fg{ag^{2}}{17.9} & \dg{8.6} & \dg{8.10} & \dg{9.5} & \dg{9.7} & \dg{10.4} & \dg{10.8} & \dg{11.5} & \dg{12.6} & \dg{14.4} & \dg{16.4} & \dg{21.3} & \dg{13.5} & \dg{d} \\
\dg{3.5} & \dg{e} &
  & & \dg{6.10} & \dg{6.14} & \dg{7.9} & \dg{7.11} & \dg{8.6} & \dg{8.8} & \dg{8.12} & \dg{9.7} & \dg{9.9} & \dg{10.6} & \dg{10.10} & \dg{11.7} & \dg{12.8} & \dg{14.6} & \dg{16.6} & \dg{21.5} & & \dg{e} \\
\dg{3.9} & \dg{f} &
  & & & \dg{6.18} & \dg{7.13} & \dg{7.15} & \dg{8.10} & \dg{8.12} & \dg{8.16} & \dg{9.11} & \dg{9.13} & \dg{10.10} & \dg{10.14} & \dg{11.11} & \dg{12.12} & \dg{14.10} & \dg{16.10} & \dg{21.9} & & \dg{f} \\
\dg{4.4} & \dg{h} &
  & & & & \dg{8.8} & \dg{8.10} & \fg{ag^{2}}{18.10} & \dg{9.7} & \dg{9.11} & \dg{10.6} & \dg{10.8} & \dg{11.5} & \dg{11.9} & \dg{12.6} & \dg{13.7} & \dg{15.5} & \dg{17.5} & \dg{22.4} & \dg{14.6} & \dg{h} \\
\dg{4.6} & \dg{i} &
  & & & & & \dg{8.12} & \dg{9.7} & \dg{9.9} & \dg{9.13} & \dg{10.8} & \dg{10.10} & \dg{11.7} & \dg{11.11} & \dg{12.8} & \dg{13.9} & \dg{15.7} & \dg{17.7} & \dg{22.6} & & \dg{i} \\
\dg{5.1} & \dg{j} &
  & & & & & & \fg{ag^{2}}{19.7} & \fg{b^{2}g}{18.8} & \dg{10.8} & \fg{ag^{2}}{20.8} & \dg{11.5} & \fg{bg}{18.4} & \dg{12.6} & \dg{13.3} & \fg{bg}{20.6} & \fg{g}{20.2} & \dg{18.2} & \dg{23.1} & \dg{15.3} & \dg{j} \\
\dg{5.3} & \dg{k} &
  & & & & & & & \dg{10.6} & \dg{10.10} & \dg{11.5} & \dg{11.7} & \dg{12.4} & \dg{12.8} & \dg{13.5} & \dg{14.6} & \dg{16.4} & \dg{18.4} & \dg{23.3} & \dg{15.5} & \dg{k} \\
\dg{5.7} & \dg{l} &
  & & & & & & & & \dg{10.14} & \dg{11.9} & \dg{11.11} & \dg{12.8} & \dg{12.12} & \dg{13.9} & \dg{14.10} & \dg{16.8} & \dg{18.8} & \dg{23.7} & & \dg{l} \\
\dg{6.2} & \dg{m} &
  & & & & & & & & & \dg{12.4} & \dg{12.6} & \dg{13.3} & \dg{13.7} & \dg{14.4} & \dg{15.5} & \dg{17.3} & \dg{19.3} & \dg{24.2} & \dg{16.4} & \dg{m} \\
\dg{6.4} & \dg{n} &
  & & & & & & & & & & \dg{12.8} & \dg{13.5} & \dg{13.9} & \dg{14.6} & \dg{15.7} & \dg{17.5} & \dg{19.5} & \dg{24.4} & & \dg{n} \\
\dg{7.1} & \dg{o} &
  & & & & & & & & & & & \dg{14.2} & \dg{14.6} & \dg{15.3} & \dg{16.4} & \dg{18.2} & \dg{20.2} & \dg{25.1} & \dg{17.3} & \dg{o} \\
\dg{7.5} & \dg{p} &
  & & & & & & & & & & & & \dg{14.10} & \dg{15.7} & \dg{16.8} & \dg{18.6} & \dg{20.6} & \dg{25.5} & & \dg{p} \\
\dg{8.2} & \dg{r} &
  & & & & & & & & & & & & & \dg{16.4} & \dg{17.5} & \dg{19.3} & \dg{21.3} & \dg{26.2} & & \dg{r} \\
\dg{9.3} & \dg{s} &
  & & & & & & & & & & & & & & \dg{18.6} & \dg{20.4} & \dg{22.4} & \dg{27.3} & \dg{19.5} & \dg{s} \\
\dg{11.1} & \dg{t} &
  & & & & & & & & & & & & & & & \dg{22.2} & \dg{24.2} & \dg{29.1} & \dg{21.3} & \dg{t} \\
\dg{13.1} & \dg{v} &
  & & & & & & & & & & & & & & & & \dg{26.2} & \dg{31.1} & & \dg{v} \\
\dg{18.0} & \dg{w} &
  & & & & & & & & & & & & & & & & & \dg{36.0} & & \dg{w} \\
\hline
\end{tabular}}
\par\smallskip
\parbox{0.95\textwidth}{\footnotesize
Each term inside this diagram is a deg-order indicating the congregate determined by an irreducible syzygy:
viz. the congregate is the product of the outside covariants in the line and column containing the deg-order, and of the literal factor (if any) placed immediately above the deg-order.
Thus, line $d$ and column $i$, $7.9$, indicates the congregate $di$, but, same line and column $j$, $17.9$ indicates the congregate $dj\cdot ag^{2}=adag^{2}$.}
\endgroup
\end{center}

% ========== p. 368 (printed p. 348) ==========

\noindent Observe as regards the foregoing diagram, that $dj^{2}$ is irreducible (since neither $dj$ nor $j^{2}$ is segregate), and similarly $j^{2}h$, $j^{4}$, \&c., are irreducible: we have thus the last or $j^{2}$ column of the diagram.

The simply divisible syzygies are infinite in number, as are also the reducible syzygies not simply divisible. There is obviously no use in writing down a simply divisible syzygy; but as regards the reducible syzygies not simply divisible, these require a calculation, and it is proper to give them as far as they have been obtained.

373. The following Table, No.\ 96, replaces Tables 88 and 89 of my Ninth Memoir. The arrangement is according to deg-orders, and the table is complete up to the deg-order 8\,.\,40: it shows for each deg-order the segregate covariants, and also the congregate covariants (if any), and the syzygies which are the expressions of these in terms of the segregates. When there are only segregates these are given in the same horizontal line with the deg-order; for instance, 5\,.\,9 $ab^{2}$, $ah$, $cd$, shows that for the deg-order 5\,.\,9 the only covariants are the segregates $ab^{2}$, $ah$, $cd$; but when there are also congregates, the segregates are arranged in the same horizontal line with the deg-order, and the congregates, each in its own horizontal line together with its expression as a linear function of the segregates: thus
\[
5\,.\,11\quad *\quad ai,\ ce \qquad bf \,\Big|\, \begin{array}{c} ai\ ce \\ -1\ +1 \end{array}
\]
the segregates are $ai$, $ce$, and there is a congregate $bf$ which is a linear function of these, $= -ai + ce$. The table gives the irreducible syzygies and also the reducible syzygies which are not simply divisible, but the simply divisible syzygies are indicated each by a reference to the divided syzygy which occurs previously in the table.

374. Any syzygy might of course be directly verified by substituting for the several covariants contained therein their expressions in terms of the coefficients and facients of the quintic. But it is to be remarked that among the syzygies, or easily deducible from them, we have (6\,.\,18) the before-mentioned equation $f^{2} = -a^{3}d + a^{2}bc - 4c^{3}$, and also a set of 17 syzygies, the left-hand sides of which are the covariants $g, h, \ldots, u, v, w$, each multiplied by $a$ or $a^{2}$, and which lead ultimately to the standard expressions of these covariants respectively, viz.\ each covariant multiplied by a proper power of $a$ can be expressed as a rational and integral function of $a, b, c, d, e, f$, linear as regards $f$. Supposing them thus expressed, a far more simple verification of any syzygy would consist in substituting therein for the several covariants their expressions in the standard form, reducing if necessary by the equation $f^{2} = -a^{3}d + a^{2}bc - 4c^{3}$: but of course, as to the syzygies used for obtaining the standard forms, this is only a verification if the standard forms have been otherwise obtained, or are assumed to be correct.

% ========== p. 369 (printed p. 349) ==========

\noindent The 17 syzygies above referred to are

\medskip

\begin{center}
\begin{tabular}{lll}
Deg-ord. & & \\
6.10 & $a^{2}g =$ & $12abd + 4b^{2}c + e^{2}$, \\
6.14 & $a^{2}h =$ & $6acd + 4bc^{2} + ef$, \\
5.11 & $ai =$ & $-bf + ce$, \\
6.6 & $aj =$ & $b^{2} - 2bh + cg + 9d^{2}$, \\
6.8 & $ak =$ & $2bi - 3de$, \\
6.12 & $al =$ & $kci - 4df$, \\
7.7 & $am =$ & $b^{2}d - cj + 3dh$, \\
7.9 & $an =$ & $b^{2}e - 6bl - 2ck - fg$, \\
8.6 & $ao =$ & $2b^{2} - ej$, \\
8.10 & $ap =$ & $-2cn - fj$, \\
9.5 & $aq =$ & $2bg - 4d^{2}f + 12dj^{2} + hj$, \\
9.7 & $ar =$ & $b^{2}k - bp + co - hi$, \\
10.8 & $as =$ & $3bdk + 3djf + 2im$, \\
12.6 & $at =$ & $vj^{2} + jp - 2mn$, \\
13.5 & $18au =$ & $2agq + b^{2}j + 6bmj - 9dt - ghj + no$, \\
14.6 & $3au =$ & $-6b^{2}d - 8b^{2}f + 2b^{2}g - 2b^{2}q + 12b^{2}h^{2} - 3et$, \\
19.5 & $18aw =$ & $3b^{2}g + b^{2}qo - 4bfo - bgmo + 18bm^{2} + 3dgjo - 18djt - 3ght - 6m^{2}o$,\\
\end{tabular}
\end{center}

\noindent the last four of these being, however, beyond the limits of the table: the expressions of $g, h, i$ are here in the standard form: the standard forms of the other covariants $j, k, \ldots, u, v, w$, will be given further on.

\medskip

\begin{center}
\textit{Table No.\ 96 (Segregates, Congregates, and Syzygies).}
\end{center}

\medskip

\begin{center}
\begin{tabular}{lll}
Deg-ord. & Congs. & Segregates. \\
\hline
1\,.\,1 & & 1 \\
3 & & \\
5 & & $a$ \\
2\,.\,2 & & $b$ \\
4 & & \\
6 & & $c$ \\
8 & & \\
10 & & $a^{2}$ \\
3\,.\,1 & & \\
3 & & $d$ \\
5 & & $e$ \\
7 & & $ab$ \\
9 & & $f$ \\
11 & & $ac$ \\
13 & & \\
15 & & $a^{3}$ \\
\end{tabular}
\end{center}

% ========== p. 370 (printed p. 350) ==========

\begin{center}
\textit{Table No.\ 96 (continued).}
\end{center}

\begin{center}
\begin{tabular}{lll}
Deg-ord. & Congs. & Segregates. \\
\hline
4\,.\,0 & & $g$ \\
2 & & $b^{2}$, $h$ \\
4 & & \\
6 & & $ad$, $bc$ \\
8 & & $ae$ \\
10 & & \\
12 & & $b$, $a^{2}b$ \\
14 & & \\
16 & & $af$ \\
18 & & \\
20 & & $a^{2}e$ \\
& & $a^{4}$ \\
\hline
5\,.\,1 & & $j$ \\
3 & & $k$ \\
5 & & $ag$, $bd$ \\
7 & & \\
9 & & $be$, $l$, $ab^{2}$, $ah$, $cd$ \\
11 & $*$ & $ai$, $ce$ \\
& $bf$ & $-1$ $+1$ \\
13 & & $a^{2}g$, $abc$ \\
15 & & \\
17 & & $a^{2}e$, $cf$ \\
19 & & \\
21 & & $a^{3}b$, $ac^{2}$ \\
23 & & \\
25 & & $a^{3}f$ \\
\hline
6\,.\,0 & & $a^{2}e$ \\
2 & & $m$, $bg$ \\
4 & & $n$ \\
6 & $*$ & $aj$, $b^{2}$, $bh$, $cg$ \\
& $d^{2}.9$ & $+1$ $-1$ $+2$ $-1$ \\
8 & $*$ & $bi$, $ak$ \\
& $de.3$ & $+1$ $+2$ \\
10 & & $c^{2}g$, $abd$, $b^{2}c$, $eh$ \\
& & $+1$ $-12$ $-4$ \quad . \\
12 & $*$ & $cdc$, $al$, $ci$ \\
& $df.3$ & $.$ $-1$ $+2$ \\
14 & $*$ & $a^{2}b^{2}$, $a^{2}h$, $acd$, $bc^{2}$ \\
& & $+1$ $-6$ $-4$ \\
16 & $*$ & $a^{2}b$, $ace$ \\
& $abf$ & \\
\end{tabular}
\end{center}

% ========== p. 371 (printed p. 351) ==========

\begin{center}
\textit{Table No.\ 96 (continued).}
\end{center}

\begin{center}
\begin{tabular}{lll}
Deg-ord. & Congs. & Segregates. \\
\hline
6\,.\,18 & $*$ & $a^{3}d$, $a^{2}bc$, $c^{3}$ \\
& $f^{2}$ & $-1$ $+1$ $-4$ \\
20 & & $a^{3}e$, $acf$ \\
22 & & $a^{4}b$, $a^{2}c^{2}$ \\
24 & & \\
26 & & $a^{3}f$ \\
28 & & \\
30 & & $a^{5}c$ \\
\hline
7\,.\,1 & & $o$ \\
3 & & $bj$, $dg$ \\
5 & $*$ & $p$, $hk$, $eg$ \\
& $dk.3$ & $0$ $+1$ $+2$ $+1$ \\
7 & & $abg$, $am$, $b^{2}d$, $ej$ \\
9 & $*$ & $an$, $bi$, $ck$, $fg$ \\
& $di.3$ & $+1$ $+6$ $+2$ $+1$ \\
& $eh$ & $+1$ $+1$ \\
11 & $*$ & $a^{2}j$, $ab^{2}$, $abh$, $acg$, $bcd$ \\
& $ad^{2}$ & $+1$ $+1$ $+4$ $+2$ $+1$ \quad S.D. 6.6, $d^{2}$ \\
& & $-1$ $+6$ $-6$ \\
13 & $*$ & $c^{2}k$, $abi$, $bce$, $el$ \\
& $cde$ & \quad\quad\quad\quad\quad\quad S.D. 6.8, $de$ \\
& $fh.3$ & $-1$ $-2$ $+3$ $-6$ \quad\quad S.D. 5.11, $bf$ \\
15 & $*$ & $a^{2}g$, $a^{2}bd$, $ab^{2}c$, $ach$, $c^{2}d$ \\
& $ae^{2}$ & \quad\quad\quad\quad\quad\quad S.D. 6.10, $e^{2}$ \\
& $fi$ & $+1$ $-1$ $+1$ $-6$ \\
17 & $*$ & $a^{2}f$, $a^{2}be$, $a^{2}d$, $c^{2}e$ \\
& $adf$ & \quad\quad\quad\quad\quad\quad S.D. 6.12, $df$ \\
& $bef$ & \quad\quad\quad\quad\quad\quad S.D. 5.11, $bf$ \\
19 & & $a^{3}b^{2}$, $a^{3}h$, $a^{2}cd$, $abc^{2}$ \\
& & \quad\quad\quad\quad\quad\quad S.D. 6.14, $ef$ \\
21 & $*$ & $a^{3}f$, $a^{3}ce$, $c^{3}f$ \\
& & $a^{4}d$, $a^{3}bc$, $ac^{3}$ \quad\quad S.D. 5.11, $bf$ \\
23 & $*$ & $a^{2}bf$ \quad\quad\quad\quad\quad\quad S.D. 6.18, $f^{2}$ \\
\end{tabular}
\end{center}

% ========== p. 372 (printed p. 352) ==========

\begin{center}
\textit{Table No.\ 96 (continued).}
\end{center}

\begin{center}
\begin{tabular}{lll}
Deg-ord. & Congs. & Segregates. \\
\hline
7\,.\,25 & & $a^{3}e$, $a^{2}cf$ \\
27 & & $a^{4}b$, $a^{2}c^{2}$ \\
29 & & $a^{3}f$ \\
31 & & \\
33 & & $a^{5}c$ \\
35 & & \\
& & $r$ \\
\hline
8\,.\,0 & & $q$ \\
2 & & $b^{2}g$, $bm$, $dj$, $gh$ \\
4 & & \\
6 & $*$ & $ao$, $bn$, $gi$ \\
& $dk.3$ & $+1$ $-3$ $-1$ \\
& $ej$ & $+1$ $-2$ \\
8 & $*$ & $a^{2}i$, $ab^{2}$, $b^{3}$, $b^{2}h$, $bcg$, $cm$ \\
& $bd^{2}$ & $-4$ $+3$ $+4$ $-6$ \quad\quad S.D. 6.6, $d^{2}$ \\
& $ek$ & $+4$ $-3$ $-4$ $+8$ $+1$ $+12$ \quad S.D. 6.8, $de$ \\
& $h^{2}.3$ & \\
10 & & $abk$, $aeg$, $ap$, $b^{2}l$, $cn$ \\
& & $+2$ \quad . \quad $+3$ $+1$ $+3$ \\
& $bhe$ & \quad $-1$ \quad . \quad . \quad $-2$ \\
& $dl.3$ & \\
& $fj$ & $+1$ \quad . \quad . \quad $+2$ $-3$ \\
& $hi.3$ & \\
12 & & $a^{2}bg$, $abn$, $ab^{2}d$, $acj$, $b^{2}c$, $bch$, $c^{2}g$ \\
& $adh$ & \quad\quad\quad\quad\quad\quad S.D. 7.7, $dh$ \\
& $be^{2}$ & \quad\quad\quad\quad\quad\quad S.D. 6.10, $e^{2}$ \\
& $cd^{2}$ & \quad\quad\quad\quad\quad\quad S.D. 6.6, $d^{2}$ \\
& $d$ & $-1$ $-2$ \quad . \quad $+1$ $-2$ $+2$ \\
& & $-1$ $+4$ $-6$ $+2$ \\
& $fk$ & $+1$ \quad . \quad . \quad $+1$ $-2$ $+1$ \\
14 & & $a^{2}n$, $abl$, $ack$, $afg$, $bci$ \\
& $adi$ & \quad\quad\quad\quad\quad\quad S.D. 7.9, $bi$ \\
& $aeh$ & \quad\quad\quad\quad\quad\quad S.D. 7.9, $dh$ \\
& $bdf$ & \quad\quad\quad\quad\quad\quad S.D. 7.9, $eh$ \\
& $cde$ & \quad\quad\quad\quad\quad\quad S.D. 6.12, $df$ \\
16 & & \quad\quad\quad\quad\quad\quad S.D. 6.8, $de$ \\
& $a^{2}d^{2}$ & \\
& $aei$, $bef$ & $a^{2}j$, $a^{2}b^{2}$, $a^{2}bh$, $a^{2}cg$, $abcd$, $b^{2}c$, $c^{2}h$ \\
& $fl.3$ & \quad\quad\quad\quad\quad\quad S.D. 6.6, $d^{2}$ \\
& & \quad\quad\quad\quad\quad\quad S.D. 7.11, $ei$ \\
& & \quad\quad\quad\quad\quad\quad S.D. 6.14, $ef$ \\
& & \quad\quad\quad\quad\quad\quad S.D. 6.10, $e^{2}$ \\
& & $+1$ $-1$ $+2$ \quad . \quad $-6$ $+6$ \\
\end{tabular}
\end{center}

% ========== p. 373 (printed p. 353) ==========

\begin{center}
\textit{Table No.\ 96 (continued).}
\end{center}

\begin{center}
\begin{tabular}{lll}
Deg-ord. & Congs. & Segregates. \\
\hline
8\,.\,18 & $*$ & $a^{2}k$, $a^{2}bi$, $abce$, $acl$, $c^{2}l$ \\
& $abde$ & \quad\quad\quad\quad\quad\quad S.D. 6.8, $de$ \\
& $abf$ & \quad\quad\quad\quad\quad\quad S.D. 5.11, $bf$ \\
& $af h$ & \quad\quad\quad\quad\quad\quad S.D. 7.13, $fh$ \\
& $cdf$ & \quad\quad\quad\quad\quad\quad S.D. 6.12, $df$ \\
20 & $*$ & $a^{2}g$, $a^{2}bd$, $a^{2}b^{2}c$, $a^{2}ch$, $ac^{2}d$, $bc$ \\
& $aby$ & \quad\quad\quad\quad\quad\quad S.D. 6.10, $e^{2}$ \\
& $afh$ & \quad\quad\quad\quad\quad\quad S.D. 7.15, $fi$ \\
& $cdf$ & \quad\quad\quad\quad\quad\quad S.D. 6.18, $f^{2}$ \\
& & \quad\quad\quad\quad\quad\quad S.D. 6.14, $ef$ \\
& $*$ & \quad\quad\quad\quad\quad\quad S.D. 6.12, $df$ \\
& $afi$ & \quad\quad\quad\quad\quad\quad S.D. 6.14, $ef$ \\
& $bf^{2}$ & \quad\quad\quad\quad\quad\quad S.D. 6.18, $f^{2}$ \\
& $cef$ & \quad\quad\quad\quad\quad\quad S.D. 5.11, $bf$ \\
22 & $*$ & $a^{2}be$, $a^{2}f$, $a^{2}ci$, $abcf$, $ac^{2}e$ \\
& & \quad\quad\quad\quad\quad\quad S.D. 6.18, $f^{2}$ \\
& $a^{2}df$ & \\
24 & $*$ & $a^{3}b$, $a^{3}b^{2}$, $a^{3}cd$, $a^{2}bc^{2}$, $c^{4}$ \\
26 & & \\
& $a^{3}kf$ & \\
28 & $*$ & $a^{4}d$, $a^{4}bc$, $a^{3}c^{2}$ \\
30 & & $a^{4}e$, $a^{3}cf$ \\
32 & & \\
34 & & $a^{5}f$ \\
36 & & $bo$, $gk$, $s$ \\
38 & & \\
40 & $*$ & $af$, $aq$, $b^{2}j$, $bdg$, $hj$ \\
\hline
9\,.\,1 & $dm.12$ & \quad $-1$ $-2$ $+1$ $+1$ \\
3 & $*$ & $ar$, $beg$, $bp$, $co$, $gl$ \\
5 & $b^{2}k.3$ & $+1$ \quad . \quad $-5$ $-1$ \quad . \quad $+1$ \\
7 & & \quad $-1$ \\
9 & $dn.3$ & \quad . \quad . \quad $-1$ $-1$ \\
& $em.3$ & \\
& $hk.3$ & $+2$ \quad . \quad $+2$ $+1$ $-1$ \\
& $j^{2}$ & $+2$ \quad . \quad $+2$ $+4$ $-1$ \\
& $*$ & . \quad $+1$ $+1$ \\
& $bdh.3$ & $a^{2}b^{2}g$, $abm$, $adj$, $agh$, $b^{2}d$, $bcj$, $cdg$ \\
& $d^{3}.27$ & \\
& $en$ & \quad $+2$ $+3$ . \quad $+1$ $+2$ $-3$ \\
& $ik$ & \quad $-1$ \quad . \quad . \quad . \quad $-3$ $+3$ \\
\end{tabular}
\end{center}

% ========== p. 374 (printed p. 354) ==========

\begin{center}
\textit{Table No.\ 96 (continued).}
\end{center}

\begin{center}
\begin{tabular}{lll}
Deg-ord. & Congs. & Segregates. \\
\hline
9\,.\,11 & $*$ & $a^{2}o$, $abn$, $agi$, $b^{2}l$, $bek$, $ceg$, $cp$ \\
& $adk$ & \quad\quad\quad\quad\quad\quad S.D. 8.6, $dk$ \\
& $ag$ & \quad\quad\quad\quad\quad\quad S.D. 8.6, $ej$ \\
& $b^{2}e$ & \quad\quad\quad\quad\quad\quad S.D. 7.9, $bi$ \\
& $bf^{2}$ & \quad . \quad $+1$ $-1$ $+6$ $+2$ $+1$ \quad S.D. 5.11, $bf$ \\
& $bdi$ & \quad . \quad . \quad $-1$ $+4$ $+2$ $+1$ \\
& $beh$ & \quad\quad\quad\quad\quad\quad S.D. 6.6, $d^{2}$ \\
& $eg$ & $+1$ $-3$ $-1$ $+2$ $+2$ \quad\quad S.D. 8.8, $ek$ \\
& $d^{2}e.9$ & $+1$ $+3$ $+1$ \quad . \quad $-3$ \quad . \quad $-3$ \quad S.D. 8.8, $h^{2}$ \\
& $fm.3$ & $+1$ $-3$ $-1$ $+2$ $+5$ \quad . \quad $-6$ \quad S.D. 7.11, $ei$ \\
13 & & $a^{2}bj$, $a^{2}dg$, $ab^{3}$, $ab^{2}h$, $abcg$, $acm$, $b^{2}cd$, $c^{3}j$ \\
& $hi.3$ & \quad\quad\quad\quad\quad\quad S.D. 7.7, $dh$ \\
& $abd^{2}$ & \\
& $aek$ & $-8$ $+4$ $-4$ $+3$ $+4$ \quad . \quad $-12$ \\
& $ah^{2}$ & $-2$ $+1$ $+2$ $-3$ $+1$ $-2$ \\
& $bei$ & \\
& $cdh$ & \\
& $d^{2}c.3$ & \\
& & \quad\quad\quad\quad\quad\quad $+2$ \\
& & $-1$ \quad . \quad . \quad $+1$ $-2$ $+1$ $-3$ $-3$ $+3$ \\
\hline
10\,.\,0 & & $bg^{2}$, $hg$, $g^{2}$, $z$, $i$ \\
2 & $*$ & $br$, $gn$, $jk$ \\
4 & $do$.6 & $+2$ \quad . \quad $-1$ \\
6 & $*$ & $aaj$, $b^{2}g$, $b^{2}m$, $bdj$, $bgh$, $cg^{2}$, $cq$ \\
& $dfg$ & \quad\quad\quad\quad\quad\quad S.D. 6.6, $d^{2}$ \\
& $eo$ & $1$ $+6$ $+1$ $+2$ $+12$ $-12$ $-2$ \\
& $hm.6$ & $-1$ $-4$ $-1$ $+4$ $+12$ $-2$ $+1$ $-3$ \\
& $bdk$ & \quad\quad\quad\quad\quad\quad $-12$ $+2$ $-1$ \\
& $bej$ & $+1$ \\
& $deg$ & $abo$, $ag^{2}$, $as$, $b^{2}n$, $bgi$, $cr$ \\
8 & $*$ & \quad\quad\quad\quad\quad\quad S.D. 8.6, $dk$ \\
& $dp$ & \quad\quad\quad\quad\quad\quad S.D. 8.6, $ej$ \\
& $hn.3$ & \quad\quad\quad\quad\quad\quad S.D. 6.8, $de$ \\
& $hi.3$ & $-5$ \quad . \quad . \quad $+3$ $+15$ $+5$ $-6$ \\
& $jl.3$ & $-6$ $+1$ \quad . \quad . \quad $+18$ $+15$ $-12$ \\
& $adm$ & \quad\quad\quad\quad\quad\quad $+13$ \quad $+1$ \\
10 & & $+1$ \quad . \quad . \quad . \quad $+3$ \\
& $b^{2}r$ & $-5$ \quad . \quad $+3$ $+15$ $+5$ $-12$ \\
& $b^{2}d^{2}.72$ & $ag^{2}$, $a^{2}q$, $ab^{2}j$, $cb^{2}dg$, $abj$, $b^{2}h$, $b^{2}cg$, $bcm$, $cdj$, $cgh$ \\
& $bek.2$ & \quad\quad\quad\quad\quad\quad $dm$\ S.D. 9.5. \\
& $bh^{2}.6$ & \\
& $d^{2}h.27$ & \\
& $eg$ & \quad\quad $+1$ $+12$ \quad . \quad $-9$ $-3$ $+4$ $+3$ $+12$ \quad . \quad $+32$ $-12$ $-4$ \\
& $g^{2}.2$ & \quad\quad $+1$ $+2$ \quad . \quad . \quad . \quad . \quad . \quad $-8$ $-32$ $+12$ $+4$ \\
& & $-1$ $-2$ \quad . \quad $+9$ $-3$ $+4$ $+3$ \quad . \quad $+4$ $+32$ $-12$ $-4$ \\
& $in.4$ & . \quad . \quad . \quad $-3$ \quad . \quad . \quad . \quad $-2$ $-8$ $+12$ $-4$ \\
& & $+1$ $-2$ \quad . \quad . \quad . \quad . \quad . \quad . \quad . \quad . \quad . \\
& $jl.4$ & . \quad . \quad . \quad $+3$ \quad . \quad . \quad . \quad $-8$ $+12$ $+1$ \\
& & \quad\quad $+3$ \\
& & \quad\quad\quad\quad\quad\quad S.D. 6.10, $e^{2}$ \\
& & . \quad . \quad . \quad . \quad . \quad $-1$ \quad . \quad $-2$ $-12$ $+6$ $+2$ \\
& & $+1$ $+2$ $-3$ $+7$ \quad . \quad . \quad . \quad $+4$ $+24$ $-24$ $-4$ \\
& & $-1$ $-2$ $+3$ $-5$ \quad . \quad . \quad . \quad . \quad . \quad . \quad . \\
& & $1$ $+2$ \quad . \quad . \quad . \quad . \quad . \quad . \quad . \quad $+12$ \\
& & . \quad . \quad . \quad $-1$ $-1$ \quad . \quad . \quad . \quad $+16$ $+12$ \\
\end{tabular}
\end{center}

% ========== p. 375 (printed p. 355) ==========

\begin{center}
\textit{Table No.\ 96 (concluded).}
\end{center}

\begin{center}
\begin{tabular}{lll}
Deg-ord. & Congs. & Segregates. \\
\hline
11\,.\,1 & & $go$, $t$ \\
3 & & $bgj$, $df$, $dq$, $jm$ \\
5 & $*$ & $b^{2}o$, $bgk$, $bs$, $cg^{2}$, $cq$, $gp$ \\
& $dr.18$ & $-2$ $+5$ $-6$ \quad . \quad $-3$ $+3$ \\
& $ho.3$ & $-2$ $+11$ $-24$ \quad . \quad $-3$ $+6$ \\
12\,.\,2 & & $jn$ \\
4 & & $km.6$ \quad $-1$ $+3$ $-1$ $+26$ \quad . \quad . \quad $+1$ $-2$ \\
13\,.\,1 & & . \quad $-3$ $+3$ \\
3 & $*$ & $gr$, $jo$ \\
& $ko$ & \\
14\,.\,0 & & $bv$, $b^{2}q$, $bgm$, $bf$, $dgj$, $g^{2}h$, $hg$ \\
2 & $*$ & . \quad $-2$ $-2$ $-4$ $+3$ \\
4 & $m\lambda.12$ & . \quad $+2$ $+1$ $+4$ $-3$ \quad . \quad $-3$ \\
& $jr.2$ & $fj$, $jq$, $b^{2}$ \\
& $nu.2$ & $bgo$, $bt$, $g^{2}k$, $gs$, $kg$ \\
& $*$ & \\
& $dgo$ & . \quad $-2$ $+1$ $-1$ . \\
& $dt.18$ & . \quad $-4$ \quad\quad $+1$ $-1$ \\
& $mr.12$ & \\
& & $b^{2}$, $bgq$, $bu$, $fm$, $gj^{2}$, $mq$, $o^{2}$ \\
& & $bgr$, $bjo$, $ghi$, $gjk$, $js$, $nq$ \\
& & \\
& & $+1$ $+2$ \quad . \quad . \quad . \quad . \quad $-1$ $+6$ $+3$ \quad . \quad . \quad . \quad . \quad . \quad . S.D. 10.4, $do$ \\
& & $+1$ $+2$ \quad . \quad . \quad . \quad . \quad $-1$ \quad . \quad $+3$ \\
\end{tabular}
\end{center}

\medskip

\begin{center}
\textit{Theory of the Canonical Form.} Art.\ Nos.\ 375 to 381, and Tables Nos.\ 97 and 98.
\end{center}

\medskip

\noindent 375. As the small italic letters have been used to represent the covariants, different letters are required for the coefficients of the quintic: using also new letters for the facients, I take the quintic to be $(a, b, c, d, e, f\,\$\,\xi, \eta)^{5}$. Considering a linear transformation of $\tfrac{1}{a}(a, b, c, d, e, f\,\$\,\xi, \eta)^{5}$, viz.
\[
\tfrac{1}{a}(a, b, c, d, e, f\,\$\,\xi - b\eta, a\eta)^{5},
\]

\end{document}
